\documentclass{beamer}
\usetheme{CambridgeUS}

\usepackage{fontspec}
\setsansfont{Fira Sans}
\setmonofont{Fira Code}
\usepackage{amsmath,amssymb,mathrsfs,wasysym}
\usepackage{unicode-math}
\unimathsetup{math-style=ISO, bold-style=ISO, mathrm=sym}
\setmathfont{Fira Math}
\setmathfont{XITS Math}[range={\mathscr,\mathbfscr}]
\setmathfont{XITS Math}[range={\mathcal,\mathbfcal},StylisticSet=1]
\usepackage{graphicx,caption,subcaption}
\usepackage{hyperref,bookmark,highlightlatex}
\hypersetup{
    colorlinks=true,
    linkcolor=blue,
    filecolor=blue,
    urlcolor=blue,
    citecolor=cyan,
}
\usepackage{geometry}
\geometry{hmargin=1.5cm}

\title{Host Galaxies of Type 1 AGNs}
\author{Zhijun Zhang}
\date{\today}

\begin{document}

\begin{frame}[plain]
    \maketitle
\end{frame}

\begin{frame}{Data Reduction and Image Decomposition}
    \begin{itemize}
        \item Data reduction including sky subtraction and PSFs.
        
        \item Decomposing and fitting the galaxy images using \texttt{GALFIT}.
        
        \item S\'{e}rsic profile:
        \[
        \mu(r) = \mu_e\mathrm{e}^{-\kappa\left[(r/R_e)^{1/n}-1\right]}
        \]
        
        \item Employ Fourier modes to characterize asymmetry, which indicates interactions or mergers.
    \end{itemize}
\end{frame}

\begin{frame}{Derived Properties}
    \begin{itemize}
        \item AGN and host galaxy luminosities and spectrum.
        
        \item Virial BH masses:
        \[
        M_{\mathrm{BH}} = \frac{fv^2 R_0}{G}
        \]
        where $R_0$ is the broad line region (BLR) in AGN. The BLR is photoionized by the central continuum source, and it is gravitationally bound to the SMBH. The normalization coefficient $f=1.31$ for $v$ set to the FWHM of $H\beta$.
        
        \item Buldge type classification using the S\'{e}rsic index and bulge-to-total light ratio.
    \end{itemize}
\end{frame}

\begin{frame}{Statistical Trends}
    \begin{itemize}
        \item Host galaxies of Radio-loud($R=f_{6\,\mathrm{cm}}/f_{4400\mathring{A}}\geqslant10$) AGNs are brighter and more bulge-dominated than their radio-quiet counterparts.
        
        \item More disk dominance, higher fraction of pseudo-bulges, higher bar frequency, and lower merger fraction implicate secular processes as the main agent for BH growth in NLS1s.
        
        \item[?] Merging features
        
    \end{itemize}
\end{frame}

\begin{frame}{Conclusion}
    \begin{itemize}
        \item H$\beta$ narrow line Type 1 AGNs prefer to host in later-type, lower-luminosity galaxies.
        
        \item The fraction of AGN hosts shows that morphological signatures of tidal interactions merger signatures is larger for more luminous hosts.
        
        \item RL AGNs generally preferentially live in earlier-type, more massive hosts, but not specially interacting or merger systems.
    \end{itemize}
\end{frame}

\begin{frame}{Discussion}
    \begin{itemize}
        \item Does dark matter distribution affect the $B/T$ measurement?
        
        \item What secular process does this article discuss?
        
        \item Why do the authors think that the trend about buldge dominating is precisely the opposite on Page 16?
    \end{itemize}
\end{frame}

\begin{frame}[plain]
    \begin{center}
        \Huge Thank you for listening!
    \end{center}
\end{frame}

\end{document}
